\section*{第 \ref{ch:b1:newton-dynamics} 章——牛顿运动定律}

\begin{solution}{exo:b1:newton-dynamics:1}
读数 $N = m(g + a)$，其中 $a$ 为向上的加速度：向上时，
$70 \times 11.81 = \qty{827}{N}$（``\qty{84}{kg}''）；向下时，$70 \times 7.81 =
\qty{547}{N}$（``\qty{56}{kg}''）；匀速时，\qty{687}{N}
（\qty{70}{kg}）；自由落体时，$0$。
\end{solution}

\begin{solution}{exo:b1:newton-dynamics:2}
两者加速度同为 $a$（绳不可伸长）：$m_1a = T$，
$m_2a = m_2g - T$，故 $a = m_2g/(m_1 + m_2) = \qty{3.3}{m/s^2}$，
$T = m_1a = \qty{6.5}{N}$。
\end{solution}

\begin{solution}{exo:b1:newton-dynamics:3}
$a = g(\sin\alpha - \mu_d\cos\alpha) = 9.81(0.50 - 0.17) = \qty{3.2}{m/s^2}$；
$v = \sqrt{2 \times 3.2 \times 5.0} = \qty{5.7}{m/s}$。与\omterm{def:b1:newton-dynamics:momentum}{质量}无关。
\end{solution}

\begin{solution}{exo:b1:newton-dynamics:4}
$k\Delta\ell = mg$：$\Delta\ell = 0.50 \times 9.81/200 = \qty{2.5}{cm}$。
设 $x$ 为相对平衡位置的位移，$m\ddot x = -kx$（重力与静态伸长
相互抵消）：$\omega_0 = \sqrt{k/m} = \qty{20}{rad/s}$，
$T = \qty{0.31}{s}$。
\end{solution}

\begin{solution}{exo:b1:newton-dynamics:5}
$\mu_s = \tan 22^\circ = 0.40$。弯道：$v \leq \sqrt{\mu_sgR}$：$\sqrt{0.8
\times 9.81 \times 50} = \qty{20}{m/s}$（\qty{71}{km/h}）；湿滑时，
\qty{14}{m/s}（\qty{50}{km/h}）。
\end{solution}

\begin{solution}{exo:b1:newton-dynamics:6}
跳伞者：$v_\infty = \sqrt{2 \times 785/(1.2 \times 0.70)} = \qty{43}{m/s}$；
$\mathrm{Re} = \rho vL/\eta \approx 1.2 \times 43 \times 0.5/1.8\times10^{-5}
\approx 10^6$：二次阻力。液滴：$m = \tfrac43\pi r^3\rho_w = \qty{4.2e-12}{kg}$，
$v_\infty = mg/6\pi\eta r = 4.1\times10^{-11}/3.4\times10^{-9} = \qty{1.2}{cm/s}$；
$\mathrm{Re} = 1.2 \times 0.012 \times 2\times10^{-5}/1.8\times10^{-5} = 0.02$：
线性阻力。
\end{solution}

\begin{solution}{exo:b1:newton-dynamics:7}
最低点：$v^2 = 2g\ell(1 - \cos 60^\circ) = g\ell$；$T = mg + mv^2/\ell =
2mg$。释放点：$v = 0$，$T = mg\cos\theta_0 = 0.5mg$。
\end{solution}

\begin{solution}{exo:b1:newton-dynamics:8}
$v = \sqrt{gR\tan\beta} = \sqrt{9.81 \times 200 \times 0.268} = \qty{23}{m/s}$
（\qty{82}{km/h}）。速度偏慢时，车有沿斜面下滑的趋势，摩擦力须沿斜面向上；
速度偏快时，摩擦力沿斜面向下，以提供
额外的\omterm{prop:b1:newton-dynamics:centripetal}{向心力}。
\end{solution}

\begin{solution}{exo:b1:newton-dynamics:9}
$T\cos\alpha = mg$，$T\sin\alpha = m\omega^2\ell\sin\alpha$：$\omega =
\sqrt{g/\ell\cos\alpha} = \sqrt{9.81/0.866} = \qty{3.4}{rad/s}$，周期
\qty{1.9}{s}；$T = mg/\cos\alpha = 1.15mg$。
\end{solution}

\begin{solution}{exo:b1:newton-dynamics:10}
$a = 27.8/8.0 = \qty{3.5}{m/s^2}$，$F = \qty{3.5}{kN}$；$F \leq \mu_smg$
给出 $\mu_s \geq a/g = 0.35$。湿路面上，$\mu_s \approx 0.4$ 将
加速度限制在约 $0.4g$，无论发动机功率多大：限制因素是抓地力，而非功率。
\end{solution}

\begin{solution}{exo:b1:newton-dynamics:11}
$v = v_\infty(1 - \eu^{-t/\tau})$，$v_\infty = g\tau = \qty{4.9}{m/s}$；
$z = \int v = v_\infty[t - \tau(1 - \eu^{-t/\tau})]$。在 \qty{2.0}{s} 时：
$4.9 \times (2.0 - 0.5 \times 0.98) = \qty{7.4}{m}$，真空中则为 $\tfrac12 gt^2
= \qty{20}{m}$。
\end{solution}

\begin{solution}{exo:b1:newton-dynamics:12}
沿径向投影（指向圆心）：$N - mg\cos\theta = mv^2/R$，故
$N = m(v^2/R + g\cos\theta)$。最高点（$\theta = \pi$）：$N = m(v^2/R - g)
\geq 0$ 当且仅当 $v \geq \sqrt{gR}$。由速率定律，$v_{\text{顶}}^2 =
v_0^2 - 4gR \geq gR$：$v_0 \geq \sqrt{5gR}$；此时在最低点，$N =
m(5g + g) = 6mg$。
\end{solution}

\begin{solution}{pb:b1:newton-dynamics:1}
\textbf{1.} $m\dot v = mg$：$v = gt$，$z = \tfrac12 gt^2$；在 \qty{10}{s} 时：
\qty{98}{m/s}，\qty{490}{m}。

\textbf{2.} $mg = \qty{785}{N}$。

\textbf{3.} 站在地面上时，地面以 $N = mg$ 向上支撑——这一
支撑力正是人所感受到的重量。自由落体时没有接触力：
重力同时作用于每一部分，任何部分都不受挤压。

\textbf{4.} 在一分钟内精确到百分之几以内，可以；忽略的因素有：
地球自转（离心项与科里奥利项，
见\cref{ch:b1:non-inertial-frames}）和风。

\textbf{5.} $F_d = \tfrac12\rho C_xSv^2 = 0.48\,v^2$（SI 单位）。

\textbf{6.} $m\,\dd v/\dd t = mg - 0.48v^2$。

\textbf{7.} $v_\infty = \sqrt{mg/0.48} = \sqrt{1635} = \qty{40}{m/s} =
\qty{146}{km/h}$。

\textbf{8.} 除以 $mg$：$\dd u/\dd t\,(v_\infty/g) = 1 - u^2$；
$\tau = 40.4/9.81 = \qty{4.1}{s}$。

\textbf{9.} $\int_0^u\dd u'/(1 - u'^2) = \operatorname{artanh}u = t/\tau$，
故 $u = \tanh(t/\tau)$。

\textbf{10.} $t = \tau\operatorname{artanh}(0.9) = 1.47\tau = \qty{6.1}{s}$；
对 $0.99$：$2.65\tau = \qty{11}{s}$。

\textbf{11.} $z = \int v_\infty\tanh(t/\tau)\dd t = v_\infty\tau\ln\cosh(t/\tau)$;
当 \qty{10}{s} 时：$166 \times \ln\cosh(2.43) = 166 \times 1.74 = \qty{290}{m}$，
与 \qty{490}{m} 相比。

\textbf{12.} $\mathrm{Re} = 1.2 \times 40 \times 0.5/1.8\times10^{-5} \approx
\num{1.3e6}$：远大于 $1000$，为平方阻力。

\textbf{13.} $v_\infty' = \sqrt{785/(0.6 \times 25)} = \sqrt{52} = \qty{7.2}{m/s}$.

\textbf{14.} 阻力 $= 15 \times 1600 = \qty{24}{kN}$；$a = (24000 - 785)/80
= \qty{290}{m/s^2} \approx 30g$：安全带会断裂，或跳跃者受伤。

\textbf{15.} 平均加速度 $a \approx (40 - 7)/3 = \qty{11}{m/s^2} \approx 1.1g$；
安全带拉力 $\approx m(g + a) \approx \qty{1.7}{kN}$。

\textbf{16.} $\tau' = 7.2/9.81 = \qty{0.73}{s}$：完全张开后的数秒之内，
下降运动已趋于稳定。

\textbf{17.} $\vect a = \vect 0$：安全带拉力 $= mg = \qty{785}{N}$。

\textbf{18.} 由平衡条件 $mg = \tfrac12\rho C_xSv^2$ 可得 $v \propto (C_xS)^{-1/2}$：
面积减半，则速率为原来的 $\sqrt2$ 倍（\qty{10}{m/s}）。

\textbf{19.} $v \propto \sqrt m$: $7.2 \times \sqrt{100/80} = \qty{8.1}{m/s}$.

\textbf{20.} $h = v^2/2g = 52/19.6 = \qty{2.6}{m}$：相当于从一层楼高处跳下，
即一层楼的高度。

\textbf{21.} $a = v^2/2d$：\qty{52}{m/s^2}（$5.3g$），作用距离 \qty{0.5}{m}，
作用力 $m(g + a) \approx \qty{5}{kN}$；若为 \qty{17}{m/s^2}（$1.8g$），
作用距离 \qty{1.5}{m}，力为 \qty{2.2}{kN}。着地时应翻滚。

\textbf{22.} 更大的阻力（以及部分升力）瞬间超过
重力，跳跃者减速，着地时速度低于 $v_\infty'$。

\textbf{23.} $m = \tfrac43\pi r^3\rho_w = \qty{4.2e-6}{kg}$, $S = \pi r^2 =
\qty{3.1e-6}{m^2}$: $v_\infty = \sqrt{2mg/\rho C_xS} = \sqrt{8.2\times10^{-5}/
1.9\times10^{-6}} = \qty{6.6}{m/s}$; $\mathrm{Re} = 1.2 \times 6.6 \times
2\times10^{-3}/1.8\times10^{-5} \approx 900$：大致属平方阻力。

\textbf{24.} $v_\infty = \qty{1.2}{cm/s}$（见\cref{exo:b1:newton-dynamics:6}）；
$\mathrm{Re} \approx 0.02 \ll 1$：斯托克斯阻力公式适用，结果自洽。

\textbf{25.} $m\dot{\vect v} = m\vect g + \vect F_{\text{阻力}}$；阻力
对小而慢的物体与 $v$ 呈线性，对大而快的物体则与速度的平方成正比。
数值如下：下落物体 \qty{40}{m/s}，伞下 \qty{7}{m/s}
（或雨滴），雾滴为 \qty{1}{cm/s}——同一规律横跨四个
\omterm{def:b1:units-dimensions:oom}{数量级}。
\end{solution}
